\par\SourceLeaf{525}Es ist mir eine um so erfreulichere Erscheinung, dass Sie mit grosser Neigung demjenigen Theile der Mathematik anhängen, der von jeher mein Lieblingsstudium gewesen ist, je seltener dieselbe ist. Ich wünsche Ihnen herzlich eine äussere Lage, wo Sie soviel als möglich Herr Ihrer Zeit und der Wahl Ihrer Arbeiten bleiben. Ich selbst wurde gleich nach dem Erscheinen meiner Disquisitiones durch andersartige Beschäftigungen, und später, durch meine äussern Verhältnisse sehr gehindert, meiner Neigung in dem Maasse nachzuhängen wie ich gewünscht hätte. Anstatt eines zweiten Theils jenes Werks, den ich früher beabsichtigte, werde ich mich aller Wahrscheinlichkeit nach darauf beschränken müssen, von Zeit zu Zeit ein Memoire über einen einzelnen Gegenstand zu liefern.
Die drei Abhandlungen dieser Art, die bisher im 16. Band der hiesigen Commentationen, und im ersten und vierten der Commentationes recentiores erschienen sind, enthalten aber (einen Theil der zweiten abgerechnet) keine von den Gegenständen, die ich schon 1801 zur Fortsetzung im Auge hatte, sondern neue; und so beziehen sich auch meine spätern Arbeiten dieser Art gleichfalls auf einen neuen Gegenstand, namentlich die Theorie der Biquadratischen Reste, die ich etwa in drei Abhandlungen zu geben denke; die erste davon wird in kurzem für den sechsten Band der Comment. rec. gedruckt werden, und die Hauptmaterialien für das Uebrige sowie für die ähnliche Theorie der cubischen Reste, ist, obgleich noch wenig davon ordentlich zu Papier gebracht ist, im Wesentlichen als abgemacht zu betrachten.
\par Empfehlen Sie mich gefälligst dem Herrn von \textsc{Humboldt}, falls er noch in Paris ist, und entschuldigen mich, dass ich jetzt nicht an ihn selbst schreibe, mit der Besorgniss, dass mein Brief ihn nicht treffen möchte, da er, wie ich höre, Paris zu verlassen die Absicht hatte.
\par\Needspace{5\baselineskip}Mit aufrichtiger Hochschätzung
\par\nobreak\noindent\begin{tabular*}{\linewidth}{@{}l@{\extracolsep{\fill}}r@{}}
&Ihr ergebenster\\
Göttingen den 13. September 1826.&C. F. \textsc{Gauss}.\\
\end{tabular*}
\BookRule{.25\textwidth}
\par\endgroup
